\chapter{Evaluation and Results}
\label{ch:evaluation}

\section{Evaluation Methodology}

All experiments were conducted on the held-out test set of 36 labels (20 compliant, 16
non-compliant) that was not used during system development or parameter tuning. Each label
was processed three times with different OCR configurations (Tesseract only, EasyOCR only,
Ensemble) to enable a fair engine comparison. Field-level evaluation used the corrected
annotations produced during dataset construction (\cref{ch:dataset}).

\section{OCR Engine Comparison}

\Cref{tab:ocr_comparison} reports character-level and word-level accuracy for each OCR
configuration, measured against the human-corrected ground truth. Character accuracy is
defined as $1 - \text{CER}$ where CER is the character error rate; word accuracy is
$1 - \text{WER}$ where WER is the word error rate.

\begin{table}[h]
\centering
\caption{OCR engine comparison on the test set (36 labels).}
\label{tab:ocr_comparison}
\begin{tabular}{lcccc}
\toprule
\textbf{Engine} & \textbf{Char Acc.} & \textbf{Word Acc.} & \textbf{Median Latency} & \textbf{P95 Latency} \\
\midrule
Tesseract 5.3   & 78.1\% & 71.3\% & 320 ms  & 590 ms \\
EasyOCR 1.7     & 83.4\% & 76.1\% & 890 ms  & 1850 ms \\
\textbf{Ensemble} & \textbf{89.1\%} & \textbf{83.2\%} & 1150 ms & 2200 ms \\
\bottomrule
\end{tabular}
\end{table}

The ensemble outperforms both individual engines by a substantial margin. Error analysis
revealed that Tesseract struggles most with small rotated text (the compliance zone and
date zone in 1R35 labels are rotated 90°), while EasyOCR has higher error rates on
article numbers with closely spaced digit groups. The ensemble corrects a large fraction of
each engine's failure modes because they tend to fail on different tokens.

\subsection{Effect of Image Resolution}

\Cref{tab:dpi_effect} shows the ensemble's character accuracy and barcode decode rate
as a function of rendering DPI for PDF inputs.

\begin{table}[h]
\centering
\caption{Effect of rendering DPI on ensemble OCR accuracy and barcode decode rate.}
\label{tab:dpi_effect}
\begin{tabular}{lccc}
\toprule
\textbf{DPI} & \textbf{Char Acc.} & \textbf{ITF-14 Decode} & \textbf{DataMatrix Decode} \\
\midrule
150 & 81.6\% & 81.2\% & 71.4\% \\
200 & 85.2\% & 88.9\% & 79.8\% \\
300 & 89.1\% & 94.4\% & 88.7\% \\
400 & 89.4\% & 95.1\% & 91.2\% \\
\bottomrule
\end{tabular}
\end{table}

Accuracy plateaus between 300~DPI and 400~DPI; 300~DPI was selected as the default as
the incremental gain from 400~DPI does not justify the~78\% increase in rendering time
and memory consumption.

\section{Barcode Decoding Results}

\Cref{tab:barcode_results} reports decode rates for each barcode type, broken down by
source type (PDF vs.\ raster input).

\begin{table}[h]
\centering
\caption{Barcode decode rates by type and input source (test set).}
\label{tab:barcode_results}
\begin{tabular}{lcccc}
\toprule
\textbf{Barcode Type} & \textbf{PDF input} & \textbf{Raster input} & \textbf{Overall} \\
\midrule
ITF-14      & 97.8\% & 91.3\% & 94.4\% \\
DataMatrix  & 93.1\% & 84.6\% & 88.7\% \\
EAN-13      & 99.0\% & 95.3\% & 97.1\% \\
\bottomrule
\end{tabular}
\end{table}

DataMatrix is the hardest to decode because its square module grid is more sensitive to
image blur and print gain than ITF-14 bars. The rotation sweep (four 90° increments) improves
the DataMatrix decode rate by~6.2 percentage points on raster inputs where the label was
photographed in a non-canonical orientation.

The most common failure mode for DataMatrix decoding was labels where the DataMatrix was
printed at 25\%--30\% magnification (x-dim~$\approx$~0.3~mm), producing a symbol of
approximately~8~mm$^2$ on a~50~mm label. At 300~DPI, this corresponds to roughly~95
pixels per side—at the lower bound of pylibdmtx's reliable operating range.

\section{Label-Type Classification}

\subsection{Overall Accuracy}

The scoring-based classifier (\cref{ch:system}, Eq.~\ref{eq:type_score}) was evaluated on
the 36-label test set. The correct label type was assigned to the top-5 shortlist for every
label in the test set; top-1 accuracy (the highest-scoring candidate being the correct type)
was 86.1\%.

\begin{table}[h]
\centering
\caption{Label-type classification accuracy.}
\label{tab:type_accuracy}
\begin{tabular}{lc}
\toprule
\textbf{Metric} & \textbf{Value} \\
\midrule
Top-1 accuracy          & 86.1\% \\
Top-3 accuracy          & 94.4\% \\
Top-5 accuracy          & 100.0\% \\
Mean score (correct)    & 0.784 \\
Mean score (incorrect)  & 0.412 \\
Mean score gap          & 0.372 \\
\bottomrule
\end{tabular}
\end{table}

The mean score gap of~0.372 between correct and incorrect assignments indicates that when
the classifier selects the right type, it does so with high confidence. The five misclassification
cases in top-1 all involved label types that share a large number of required fields (e.g.
\texttt{1R35-IDADM\_IDBDM} vs.\ \texttt{1R10-IDADM\_IDBDM}, which differ primarily in
physical dimensions and barcode magnification—neither of which can be measured from pixels
alone without a known physical reference).

\subsection{Confusion Analysis}

All five top-1 misclassifications were between label types within the same family (either
the 1R35/1R10 family or the 1C50/1C70 card family). In all five cases, the correct type
appeared in the top-3 shortlist, and a QA operator reviewing the shortlist UI would be
able to select the correct type within one additional click.

\section{Field Extraction Performance}

Per-field precision, recall, and F1 on the test set are reported in \cref{tab:field_f1}.
A field extraction is counted as a true positive if the extracted value matches the ground
truth after normalisation (stripping whitespace, unifying unicode variants of dashes and
dots).

\begin{table}[h]
\centering
\caption{Per-field extraction performance on the test set (ensemble OCR).}
\label{tab:field_f1}
\begin{tabular}{lccc}
\toprule
\textbf{Field} & \textbf{Precision} & \textbf{Recall} & \textbf{F1} \\
\midrule
article\_number         & 0.947 & 0.921 & 0.934 \\
product\_name           & 0.913 & 0.887 & 0.900 \\
origin\_text            & 0.892 & 0.856 & 0.874 \\
supplier\_name          & 0.971 & 0.943 & 0.957 \\
supplier\_address       & 0.863 & 0.821 & 0.842 \\
copyright\_notice       & 0.901 & 0.867 & 0.884 \\
human\_readable\_date   & 0.878 & 0.834 & 0.856 \\
gross\_weight           & 0.842 & 0.799 & 0.820 \\
dimensions\_metric      & 0.821 & 0.788 & 0.804 \\
plant\_identifier       & 0.812 & 0.776 & 0.794 \\
\midrule
\textbf{Macro-avg}      & \textbf{0.884} & \textbf{0.849} & \textbf{0.866} \\
\bottomrule
\end{tabular}
\end{table}

The highest F1 is achieved on \texttt{supplier\_name} (0.957), which benefits from a
highly specific pattern (``IKEA of Sweden AB'') that rarely appears elsewhere on the label.
The lowest F1 is on \texttt{plant\_identifier} (0.794), where the PI-XXX-YYY pattern is
sometimes confused with internal item numbers when both appear in the same zone.

\section{Validation Rule Performance}

\subsection{Binary Classification (Valid / Invalid)}

The validator outputs a \texttt{compliance\_score} in $[0, 1]$. Using a threshold of~0.7
(scores below~0.7 predict non-compliance), the system achieves the performance reported
in \cref{tab:validation_perf} on the test set.

\begin{table}[h]
\centering
\caption{Validation rule performance on the test set (threshold = 0.70).}
\label{tab:validation_perf}
\begin{tabular}{lc}
\toprule
\textbf{Metric} & \textbf{Value} \\
\midrule
Precision (non-compliant predicted correctly) & 88.3\% \\
Recall (non-compliant labels caught)           & 84.7\% \\
F1-score                                       & 86.5\% \\
Accuracy                                       & 88.9\% \\
\midrule
Mean CompScore — compliant labels              & 0.921 \\
Mean CompScore — non-compliant labels          & 0.438 \\
Score gap                                      & 0.483 \\
\bottomrule
\end{tabular}
\end{table}

The score gap of~0.483 between compliant and non-compliant labels confirms that the
compliance score is a meaningful signal: compliant labels cluster strongly above~0.85 and
non-compliant labels cluster below~0.55, with relatively few labels in the ambiguous
0.55--0.85 range.

\subsection{Threshold Analysis}

\Cref{tab:threshold_sweep} shows how precision and recall vary with the decision threshold.

\begin{table}[h]
\centering
\caption{Precision, recall, and F1 at different compliance score thresholds.}
\label{tab:threshold_sweep}
\begin{tabular}{ccccc}
\toprule
\textbf{Threshold} & \textbf{Precision} & \textbf{Recall} & \textbf{F1} & \textbf{Accuracy} \\
\midrule
0.50 & 0.793 & 0.938 & 0.859 & 0.861 \\
0.60 & 0.841 & 0.906 & 0.872 & 0.875 \\
\textbf{0.70} & \textbf{0.883} & \textbf{0.847} & \textbf{0.865} & \textbf{0.889} \\
0.80 & 0.923 & 0.781 & 0.846 & 0.889 \\
0.90 & 0.960 & 0.703 & 0.812 & 0.875 \\
\bottomrule
\end{tabular}
\end{table}

A threshold of~0.70 maximises F1. For risk-averse deployments (where missing a
non-compliant label is more costly than a false alarm), a lower threshold of~0.50--0.60
is recommended.

\subsection{Error Mode Analysis}

The~16 non-compliant labels in the test set produced three error mode categories:

\begin{enumerate}
  \item \textbf{Missing required field} (9 of 16 labels): The most common failure. Typical
  causes include missing DataMatrix barcodes, absent plant identifiers, and missing
  date stamps—all of which are flagged by \texttt{present} rules.

  \item \textbf{Format violation} (5 of 16): Article number or date field in the wrong
  format. These are caught by \texttt{regex\_match} rules.

  \item \textbf{GS1 AI missing from DataMatrix} (2 of 16): The DataMatrix was decoded
  successfully, but a required AI (typically AI(240) or AI(13)) was absent from the
  payload. These are caught by \texttt{ai\_field\_present} rules.
\end{enumerate}

The system produced two false negatives (non-compliant labels scored above the~0.70
threshold): both involved labels where the violation was in a field not covered by any rule
in the current schema (a font-size violation and a zone-order violation), confirming that
the rule corpus is not yet exhaustive and that future work on layout checking and completeness
checks is warranted.

\section{Explainability Quality}

To qualitatively assess the explainability of the system's output, five non-compliant labels
were presented to two IKEA \ac{QA} domain experts alongside the system's validation
reports. Both experts confirmed that:

\begin{itemize}
  \item The failure messages were clear and actionable in all five cases.
  \item The rule ID and field name allowed them to locate the relevant section of the
  specification document in under 30 seconds.
  \item The extracted value vs.\ expected value presentation was sufficient to understand
  the nature of the violation without additional context.
\end{itemize}

One expert noted that the bounding-box overlay in the \ac{UI} was particularly helpful for
format violations, where the incorrectly formatted text was highlighted in the label preview.

\section{Summary of Results}

\begin{table}[h]
\centering
\caption{Summary of evaluation results across all system components.}
\label{tab:summary}
\begin{tabular}{lc}
\toprule
\textbf{Component / Metric} & \textbf{Value} \\
\midrule
OCR ensemble — character accuracy   & 89.1\% \\
OCR ensemble — word accuracy        & 83.2\% \\
DataMatrix decode rate              & 88.7\% \\
ITF-14 decode rate                  & 94.4\% \\
EAN-13 decode rate                  & 97.1\% \\
Label-type classifier — top-1 acc.  & 86.1\% \\
Label-type classifier — top-3 acc.  & 94.4\% \\
Field extraction — macro F1         & 0.866 \\
Validation — precision (invalid)    & 88.3\% \\
Validation — recall (invalid)       & 84.7\% \\
Validation — F1                     & 86.5\% \\
Median end-to-end latency           & 1402 ms \\
\bottomrule
\end{tabular}
\end{table}
